% !TEX root = ../main.tex

\chapter{Problem Statement}
\label{cha:problem}

Programming, as it is taught and practiced, excludes those who do not fit its dominant mould.
This exclusion breaks down into three interrelated tensions that are worth examining separately.

\section{Interface Accessibility}
\label{sec:problem:accessibility}
The keyboard and the screen are historically situated artefacts.
They were designed for users with binaural manual dexterity, classical vision and sustained fine motor control.
Any deviation from that profile, such as the temporary or permanent loss of mobility in an arm, dyslexia, low vision, early childhood, or neurodivergence, turns the interface into an active barrier~\cite{ko23,oleson23}.
While disciplines such as music or painting have developed multiple modes of access over centuries, programming has clung to a single channel which, even though it offers accessibility solutions, is not intrinsically inclusive.

\section{Cognitive Entry Threshold}
\label{sec:problem:cognitive}
Learning to program requires mastering an exotic syntax, an abstract semantics and a mental model of execution that the beginner cannot directly inspect.
Systematic studies of the learning process~\cite{bosse17,robins03} document consistent patterns of difficulty across institutions.
The consequence is a \emph{filter} that operates before the learner can even evaluate whether programming is of interest to them.

\section{Paradigmatic Monotony}
\label{sec:problem:monotony}
Most contemporary languages descend, with surface variations, from the same syntactic lineage.
The proposals that challenged this precept, such as visual, block-based or end-user programming languages, improved textual accessibility but, for the most part, preserved the fundamental assumption that the program is inscribed on a flat two-dimensional surface.
Carrying programming into three-dimensional space, into touch and into the manipulation of objects, remains under-explored, despite robust evidence on the benefits of manipulative learning in other disciplines~\cite{montessori12,marshall07}.

\section{Research Question}
\label{sec:problem:question}
From the three tensions above we formulate the guiding question:
\emph{is it possible to design a tangible, Turing-complete and aesthetically open programming language that serves as an entry point to computation for more diverse audiences, without sacrificing the formal rigour on which any conventional language depends?}
The rest of this document develops an affirmative answer while taking different perspectives and alternatives into account.

\endinput
